\documentclass[10pt,a4paper]{amsart}
\usepackage[margin=19mm]{geometry}
\usepackage{amsmath,amssymb,mathtools}
\usepackage[scr=rsfs,cal=euler]{mathalfa}
\usepackage{xfrac}
\usepackage[unicode,hidelinks,bookmarks=false]{hyperref}
\newcommand{\ie}{i.e.\ }
\newcommand{\cf}{cf.\ }
\newcommand{\resp}{respectively\ }
\newcommand{\Subsection}{\subsection}
\providecommand{\nobreakdash}{\nobreak\hbox{-}}
\setlength{\emergencystretch}{2em}
\multlinegap0pt
\usepackage{bm}
\newtheorem{proposition}{Proposition}
\newtheorem{lem}{Lemma}
\newcommand{\Set}{\operatorname{Set}}
\newcommand{\id}{\operatorname{id}} % identity map
\newcommand{\ive}[1]{\left[ #1 \right]}
\newcommand{\tup}[1]{\left( #1 \right)}
\title{The left-to-right minima basis of the group algebra of the symmetric group (updated version)}
\author{Darij Grinberg, Ekaterina A. Vassilieva}
\date{Source: arXiv:2601.02952v2 (CC BY 4.0)}
\begin{document}
\maketitle

\noindent\emph{Source and licence.} The left-to-right minima basis of the group algebra of the symmetric group (updated version). Version arXiv:2601.02952v2 (CC BY 4.0), distributed by arXiv under the Creative Commons Attribution 4.0 International licence, http://creativecommons.org/licenses/by/4.0/, as recorded in the arXiv licence metadata for this exact version and shown on the abstract page https://arxiv.org/abs/2601.02952v2. Full licence notice: \texttt{licenses/arxiv-2601.02952-CC-BY-4.0.txt}.\par\smallskip

\noindent\textit{Editorial scope.} Counted unit: the article's lem lem.wuv (source0.tex lines 1283--1309). The article's complete original proof of it is delivered with it (source0.tex lines 1311--1491, one \texttt{proof} environment of 181 lines). No mathematics is written, repaired or re-derived: the statement and the proof are the article's own lines, in the article's own notation, and no retained line is substituted or re-flowed.  Retained uncounted as the source-local context the proof needs: lines 280--285; lines 964--973.  Nothing was omitted from any retained span, so the package carries no omission record.  No retained line contains a citation, so the wrapper carries no bibliography and no bibliography entry.  No retained line contains a figure environment, an image-inclusion command or drawing code, so no image is delivered and the package declares no asset dependency.  Editorial formatting only, in addition: an AMS \texttt{amsart} preamble that restates the article's own declarations and macros used by the retained lines, namely the environments \texttt{proposition}, \texttt{lem} on separate counters, together with the standard packages those lines need.  The retained environments print in this document as Proposition 1, Lemma 1 in order of appearance, whereas the article prints the same items with the numbering of its own full document.  The complete native source of the article as distributed in the arXiv e-print for this exact version is bundled unchanged as \texttt{sources/192199/source0.tex}.
\begin{align}
\Set(\alpha)_i := \ive{ \alpha_{\le i-1} + 1, \ \alpha_{\le i} }
\text{ of } \ive{n}
\qquad \text{for all } i \in [k].
\label{eq.alpha-blocks}
\end{align}

\begin{proposition}
\label{prop.dynkin.Vsh}
Let $S$ be a nonempty subset of $\left[  n\right]  $.
Then,
\[
V^{S}=\sum\left(  -1\right)  ^{\operatorname*{val}w-1}w,
\]
where the sum ranges over all words $w$ that contain each element of $S$ exactly
once (and no other letters) and have V-shape.
\end{proposition}

\begin{lem}
\label{lem.wuv}Let $w,u\in S_{n}$. Assume that the compositions $\beta=\left(
\beta_{1},\beta_{2},\ldots,\beta_{p}\right)  :=\operatorname{cLRM}^{\prime
}(w)$ and $\gamma=\left(  \gamma_{1},\gamma_{2},\ldots,\gamma_{p}\right)
:=\operatorname{cLRM}^{\prime}(u)$ are anagrams of one another, and let
$\chi\in S_{p}$ be a permutation satisfying%
\begin{equation}
\beta_{\chi\left(  i\right)  }=\gamma_{i}\qquad\text{for each }i\in\left[
p\right]  .
\label{eq.lem.wuv.bg}
\end{equation}

Let $v\in S_{n}$ be a permutation such that the word $\underline{v}\in
W_{n,n}$ appears in the expansion of the product $V^{\operatorname{Set}%
(\beta)_{\chi(1)}}\cdots V^{\operatorname{Set}(\beta)_{\chi(p)}}$.
Equivalently, let
$v\in S_{n}$ be a permutation that sends each $\operatorname*{Set}\left(
\gamma\right)  _{i}$ to $\operatorname*{Set}\left(  \beta\right)
_{\chi\left(  i\right)  }$ and has V-shape on each $\operatorname*{Set}\left(
\gamma\right)  _{i}$
(this is equivalent by Proposition~\ref{prop.dynkin.Vsh} and by
\eqref{eq.lem.wuv.bg}). Assume that
\[
vu=w.
\]
Then, $\chi= \id $ and $v= \id $ and $\gamma=\beta$.
\end{lem}

\begin{proof}
For each $i\in\left[  p\right]  $, we define a claim $\mathcal{C}\tup{i}$
as follows:
\begin{quote}
\textit{Claim $\mathcal{C}\tup{i}$:}
We have $\chi\tup{i} = i$
and $\gamma_i = \beta_i$
and $\Set\tup{\beta}_i = \Set\tup{\gamma}_i$
and $\left. v\mid_{\Set\tup{\beta}_i} \right. = \id$.
\end{quote}
% denote the claim saying that
% \begin{align*}
% \chi\left(  i\right)  =i\ \ \ \ \ \ \ \ \ \   & \text{and}%
% \ \ \ \ \ \ \ \ \ \ \gamma_{i}=\beta_{i}\ \ \ \ \ \ \ \ \ \ \text{and}%
% \ \ \ \ \ \ \ \ \ \ \operatorname*{Set}\left(  \beta\right)  _{i}%
% =\operatorname*{Set}\left(  \gamma\right)  _{i}\\
% & \text{and}\ \ \ \ \ \ \ \ \ \ v\mid_{\operatorname*{Set}\left(
% \beta\right)  _{i}}= \id .
% \end{align*}
It clearly suffices to show that $\mathcal{C}\left(  i\right)  $ holds for all
$i\in\left[  p\right]  $. We will prove this by strong descending induction
on $i$. That is, we fix $i\in\left[  p\right]  $, and we assume that
all the ``higher'' claims
$\mathcal{C}\left(  i+1\right)  ,\mathcal{C}\left(  i+2\right)  ,\ldots
,\mathcal{C}\left(  p\right)  $ hold. We must then prove $\mathcal{C}\left(
i\right)  $.

By \eqref{eq.alpha-blocks},
we have $\operatorname*{Set}\left(
\beta\right)  _{i}=\left[  \beta_{\leq i-1}+1,\ \beta_{\leq i}\right]  $
and $\operatorname*{Set}\left(  \gamma\right)  _{i}=\left[
\gamma_{\leq i-1}+1,\ \gamma_{\leq i}\right]  $.

If $\sigma\in S_{n}$ is a permutation and $j\leq\left\vert \operatorname*{LRM}%
\sigma\right\vert $ is a positive integer, then we let $\operatorname*{lrm}%
\nolimits_{j}\sigma$ denote the $j$-th largest LRM of $\sigma$. We extend this
to $j=0$ by setting $\operatorname*{lrm}\nolimits_{0}\sigma:=n+1$. We note
that the LRMs of a permutation $\sigma$ appear in decreasing order when
reading $\sigma$ from left to right, so that $\operatorname*{lrm}%
\nolimits_{j}\sigma$ is the $j$-th LRM of $\sigma$ encountered in such a read.
Thus, it is easy to see that each $\sigma\in S_{n}$ and each positive
$j\leq\left\vert \operatorname*{LRM}\sigma\right\vert $ satisfy%
\begin{equation}
\operatorname*{lrm}\nolimits_{j}\sigma=\sigma\left(  \min\left(  \sigma
^{-1}\left(  \left[  \operatorname*{lrm}\nolimits_{j-1}\sigma-1\right]
\right)  \right)  \right)  .\label{pf.lem.wuv.lrm-min}%
\end{equation}
(Indeed, this is just saying that if we remove all numbers larger than
$\operatorname*{lrm}\nolimits_{j-1}\sigma-1$ from the one-line notation of
$\sigma$, then the first of the remaining numbers will be $\operatorname*{lrm}%
\nolimits_{j}\sigma$; but this is clear from the definition of an LRM.)

But $\beta=\operatorname{cLRM}^{\prime}\left(  w\right)  $ shows that
$\beta_{\leq i-1}+1$ is the $i$-th smallest LRM of $w$, and hence is the
$\left(  p+1-i\right)  $-th largest LRM of $w$ (since $w$ has $\ell\left(
\operatorname{cLRM}^{\prime}(w)\right)  =\ell\tup{\beta} = p$
many LRMs in total). That is,
$\beta_{\leq i-1}+1=\operatorname*{lrm}\nolimits_{p+1-i}w$. Likewise,
$\beta_{\leq i}+1=\operatorname*{lrm}\nolimits_{p-i}w$. But
(\ref{pf.lem.wuv.lrm-min}) (applied to $\sigma=w$ and $j=p+1-i$) yields
\[
\operatorname*{lrm}\nolimits_{p+1-i}w=w\left(  \min\left(  w^{-1}\left(
\left[  \operatorname*{lrm}\nolimits_{p-i}w-1\right]  \right)  \right)
\right)  .
\]
Using $\beta_{\leq i-1}+1=\operatorname*{lrm}\nolimits_{p+1-i}w$ and
$\beta_{\leq i}+1=\operatorname*{lrm}\nolimits_{p-i}w$, this rewrites as%
\begin{equation}
\beta_{\leq i-1}+1=w\left(  \min\left(  w^{-1}\left(  \left[  \beta_{\leq
i}\right]  \right)  \right)  \right)  .\label{pf.lem.wuv.4b}%
\end{equation}
Similarly,%
\begin{equation}
\gamma_{\leq i-1}+1=u\left(  \min\left(  u^{-1}\left(  \left[  \gamma_{\leq
i}\right]  \right)  \right)  \right)  .\label{pf.lem.wuv.4g}%
\end{equation}

However, our induction hypothesis yields $\beta_{j}=\gamma_{j}$ for all $j>i$,
and thus easily $\beta_{\leq i}=\gamma_{\leq i}$ (since $\left\vert
\beta\right\vert =n=\left\vert \gamma\right\vert $). Furthermore, it shows
that $\left. v\mid_{\operatorname*{Set}\left(  \beta\right)  _{j}} \right.
= \id $ for all $j>i$, and thus
$\left. v\mid_{\left[  n\right]  -\left[  \beta_{\leq
i}\right]  }\right. = \id $,
so that
$v\tup{\left[  n\right]  -\left[  \beta_{\leq
i}\right]} = \left[  n\right]  -\left[  \beta_{\leq
i}\right]$.
By taking complements, we obtain
$v\tup{\left[  \beta_{\leq i}\right]} = \left[  \beta_{\leq i}\right]$,
thus
$v^{-1}\tup{\left[  \beta_{\leq i}\right]} = \left[  \beta_{\leq i}\right]$.
But $vu = w$, so that $w^{-1}=u^{-1}v^{-1}$ and therefore
$w^{-1}\left( \left[  \beta_{\leq i}\right]  \right)
= u^{-1}\left( v^{-1} \left( \left[  \beta_{\leq i}\right]  \right)\right)
= u^{-1}\left(  \left[  \beta_{\leq i}\right]  \right)  $,
because
$v^{-1}\tup{\left[  \beta_{\leq i}\right]} = \left[  \beta_{\leq i}\right]$.
% , so that $\left.
% v^{-1}\mid_{\left[  n\right] -\left[  \beta_{\leq i}\right]  } \right.
% = \id $ and therefore
% $w^{-1}\mid_{\left[  n\right]  -\left[  \beta_{\leq i}\right]  }=u^{-1}%
% \mid_{\left[  n\right]  -\left[  \beta_{\leq i}\right]  }$ (since $vu=w$
% entails $w^{-1}=u^{-1}v^{-1}$). Thus, $w^{-1}\left(  \left[  n\right]
% -\left[  \beta_{\leq i}\right]  \right)  =u^{-1}\left(  \left[  n\right]
% -\left[  \beta_{\leq i}\right]  \right)  $ and therefore $w^{-1}\left(
% \left[  \beta_{\leq i}\right]  \right)  =u^{-1}\left(  \left[  \gamma_{\leq
% i}\right]  \right)  $ as well.

Now, applying $v$ to (\ref{pf.lem.wuv.4g}), we find
\begin{align}
v\left(  \gamma_{\leq i-1}+1\right)    & =v\left(  u\left(  \min\left(
u^{-1}\left(  \left[  \gamma_{\leq i}\right]  \right)  \right)  \right)
\right)  \nonumber\\
& =w\left(  \min\left(  u^{-1}\left(  \left[  \beta_{\leq i}\right]  \right)
\right)  \right)  \ \ \ \ \ \ \ \ \ \ \left(  \text{since }vu=w
\text{ and } \gamma_{\leq i} = \beta_{\leq i} \right)
\nonumber\\
& =w\left(  \min\left(  w^{-1}\left(  \left[  \beta_{\leq i}\right]  \right)
\right)  \right)  \ \ \ \ \ \ \ \ \ \ \left(  \text{since }w^{-1}\left(
\left[  \beta_{\leq i}\right]  \right)  =u^{-1}\left(  \left[  \beta_{\leq
i}\right]  \right)  \right)  \nonumber\\
& =\beta_{\leq i-1}+1\ \ \ \ \ \ \ \ \ \ \left(  \text{by (\ref{pf.lem.wuv.4b}%
)}\right)  \label{pf.lem.wuv.v1}\\
& \in\left[  \beta_{\leq i-1}+1,\ \beta_{\leq i}\right]  =\operatorname*{Set}%
\left(  \beta\right)  _{i}.\label{pf.lem.wuv.vSet}%
\end{align}
But $\gamma_{\leq i-1}+1\in\left[  \gamma_{\leq i-1}+1,\ \gamma_{\leq
i}\right]  =\operatorname*{Set}\left(  \gamma\right)  _{i}$, so that%
\[
v\left(  \gamma_{\leq i-1}+1\right)  \in v\left(  \operatorname*{Set}\left(
\gamma\right)  _{i}\right)  =\operatorname*{Set}\left(  \beta\right)
_{\chi\left(  i\right)  }%
\]
(by the requirements on $v$). Confronting this with (\ref{pf.lem.wuv.vSet}),
we conclude that $\chi\left(  i\right)  =i$, since otherwise
$\operatorname*{Set}\left(  \beta\right)  _{\chi\left(  i\right)  }$ would be
disjoint from $\operatorname*{Set}\left(  \beta\right)  _{i}$. Thus,
(\ref{eq.lem.wuv.bg}) becomes $\gamma_{i}=\beta_{i}$. Thus, $\gamma_{\leq
i-1}=\beta_{\leq i-1}$ (by subtracting $\gamma_{i}=\beta_{i}$ from
$\gamma_{\leq i}=\beta_{\leq i}$). Combining this with $\gamma_{\leq i}%
=\beta_{\leq i}$, we conclude that $\operatorname*{Set}\left(  \beta\right)
_{i}=\operatorname*{Set}\left(  \gamma\right)  _{i}$.

Recall that the permutation $v$ has V-shape on $\operatorname*{Set}\left(
\gamma\right)  _{i} =\left[
\gamma_{\leq i-1}+1,\ \gamma_{\leq i}\right] $. That is, the word
$\tup{v\tup{\gamma_{\leq i-1} + 1}, v\tup{\gamma_{\leq i-1} + 2}, \ldots,
v\tup{\gamma_{\leq i}}}$ has V-shape.
However, the letters of this word are $\gamma_{\leq i-1} + 1,
\gamma_{\leq i-1} + 2, \ldots, \gamma_{\leq i}$ in some order
(since $v$ sends $\operatorname*{Set}\left(
\gamma\right)  _{i}$ to $\operatorname*{Set}\left(  \beta\right)
_{\chi\left(  i\right)  }=\operatorname*{Set}\left(  \beta\right)
_{i}=\operatorname*{Set}\left(  \gamma\right)  _{i}$),
and moreover the \textbf{leftmost} of these letters is the
\textbf{smallest} of them (since $\gamma_{\leq
i-1}=\beta_{\leq i-1}$ allows us to rewrite (\ref{pf.lem.wuv.v1}) as $v\left(
\gamma_{\leq i-1}+1\right)  =\gamma_{\leq i-1}+1$).
But a word that has V-shape and starts with its smallest letter
must be increasing; thus, in particular, the word
$\tup{v\tup{\gamma_{\leq i-1} + 1}, v\tup{\gamma_{\leq i-1} + 2}, \ldots,
v\tup{\gamma_{\leq i}}}$ is increasing.
Since its letters are $\gamma_{\leq i-1} + 1,
\gamma_{\leq i-1} + 2, \ldots, \gamma_{\leq i}$, we thus obtain
\[
\tup{v\tup{\gamma_{\leq i-1} + 1}, v\tup{\gamma_{\leq i-1} + 2}, \ldots,
v\tup{\gamma_{\leq i}}}
= \tup{\gamma_{\leq i-1} + 1,
\gamma_{\leq i-1} + 2, \ldots, \gamma_{\leq i}}.
\]
In other words,
$\left. v\mid_{\operatorname*{Set}\left(  \gamma\right)  _{i}}\right.
= \id $.
That is, $\left. v\mid_{\operatorname*{Set}\left(  \beta\right)  _{i}}%
\right. = \id $ (since $\operatorname*{Set}\left(  \beta\right)
_{i}=\operatorname*{Set}\left(  \gamma\right)  _{i}$).
We have now proved all four parts of our claim $\mathcal{C}\left(  i\right)
$. Thus, the induction step is complete, and as we said, this finishes the
proof of the lemma.
\end{proof}

\end{document}
